%% Nothing to modify here.
%% make sure to include this before anything else
% \renewcommand*\oldstylenums[1]{{\firaoldstyle #1}}
\documentclass[10pt,aspectratio=169]{beamer}
\usetheme{metropolis}
\usenavigationsymbolstemplate{}
\usepackage{tgcursor}
\usepackage{appendixnumberbeamer}
\usepackage{xcolor}
\usepackage[scale=2]{ccicons}
\usepackage[sfdefault]{FiraSans} %% option 'sfdefault' activates Fira Sans as the default text font
\usepackage[T1]{fontenc}
\usepackage{pgfplots}
\usepgfplotslibrary{dateplot}

\usepackage{xspace}
\newcommand{\themename}{\textbf{\textsc{metropolis}}\xspace}
% packages
\usepackage{color}
\usepackage{listings}
\usepackage{multicol}

% color definitions
\definecolor{mygreen}{rgb}{0,0.6,0}
\definecolor{mygray}{rgb}{0.5,0.5,0.5}
\definecolor{mymauve}{rgb}{0.58,0,0.82}

% re-format the title frame page
\makeatletter
\def\supertitle#1{\gdef\@supertitle{#1}}%
% \setbeamertemplate{title page}
% {
%     \vbox{}
%     \vfill
%     \begin{centering}
%         \begin{beamercolorbox}[sep=8pt,center]{title}
%             \usebeamerfont{supertitle}\@supertitle
%         \end{beamercolorbox}
%         \begin{beamercolorbox}[sep=8pt,center]{title}
%             \usebeamerfont{title}\inserttitle\par%
%             \ifx\insertsubtitle\@empty%
%             \else%
%                 \vskip0.25em%
%                 {\usebeamerfont{subtitle}\usebeamercolor[fg]{subtitle}\insertsubtitle\par}%
%             \fi%
%         \end{beamercolorbox}%
%         \vskip1em\par
%         \begin{beamercolorbox}[sep=8pt,center]{author}
%             \usebeamerfont{author}\insertauthor
%         \end{beamercolorbox}
%         \begin{beamercolorbox}[sep=8pt,center]{institute}
%             \usebeamerfont{institute}\insertinstitute
%         \end{beamercolorbox}
%         \begin{beamercolorbox}[sep=8pt,center]{date}
%             \usebeamerfont{date}\insertdate
%         \end{beamercolorbox}\vskip0.5em
%         {\usebeamercolor[fg]{titlegraphic}\inserttitlegraphic\par}
%     \end{centering}
%     \vfill
% }
\makeatother

% insert frame number
\expandafter\def\expandafter\insertshorttitle\expandafter{%
    \insertshorttitle\hfill%
    \insertframenumber\,/\,\inserttotalframenumber}

% preset-listing options
\lstset{
    backgroundcolor=\color{white},
    % choose the background color;
    % you must add \usepackage{color} or \usepackage{xcolor}
    basicstyle=\footnotesize,
    % the size of the fonts that are used for the code
    breakatwhitespace=false,
    % sets if automatic breaks should only happen at whitespace
    breaklines=true,                 % sets automatic line breaking
    captionpos=b,                    % sets the caption-position to bottom
    commentstyle=\color{mygreen},    % comment style
    % deletekeywords={...},
    % if you want to delete keywords from the given language
    extendedchars=true,
    % lets you use non-ASCII characters;
    % for 8-bits encodings only, does not work with UTF-8
    frame=single,                    % adds a frame around the code
    keepspaces=true,
    % keeps spaces in text,
    % useful for keeping indentation of code
    % (possibly needs columns=flexible)
    keywordstyle=\color{blue},       % keyword style
    % morekeywords={*,...},
    % if you want to add more keywords to the set
    numbers=left,
    % where to put the line-numbers; possible values are (none, left, right)
    numbersep=5pt,
    % how far the line-numbers are from the code
    numberstyle=\tiny\color{mygray},
    % the style that is used for the line-numbers
    rulecolor=\color{black},
    % if not set, the frame-color may be changed on line-breaks
    % within not-black text (e.g. comments (green here))
    stepnumber=1,
    % the step between two line-numbers.
    % If it's 1, each line will be numbered
    stringstyle=\color{mymauve},     % string literal style
    tabsize=4,                       % sets default tabsize to 4 spaces
    title=\lstname
    % show the filename of files included with \lstinputlisting;
    % also try caption instead of title
}

% macro for code inclusion
\newcommand{\includecode}[2][c]{
    \lstinputlisting[caption=#2, style=custom#1]{#2}
}
	% nothing to do here
\usepackage[german]{babel}

\usepackage[utf8]{inputenc}



\lstset{
  language = Java,
  showspaces = false,
  showtabs = false,
  showstringspaces = false,
  escapechar = @,
  belowskip=-1.5em
}

\def\ContinueLineNumber{\lstset{firstnumber=last}}
\def\StartLineAt#1{\lstset{firstnumber=#1}}
\let\numberLineAt\StartLineAt

\makeatletter
 % TODO modify this if you have not already done so


% meta-information
\usepackage{tikz}
\usepackage{hyperref}
\hypersetup{
	colorlinks=true,
	linkcolor=darkgray,
	urlcolor=blue,
}


\title{Java}
\subtitle{Object Oriented  Programming}
\date{\today}
\author{Florian Kluge, Moritz Schulz}
\institute{Florian.Kluge@mailbox.tu-dresden.de, Moritz.Schulz2@mailbox.tu-dresden.de}
% \titlegraphic{\hfill\includegraphics[height=1.5cm]{logo.pdf}}

\begin{document}
\maketitle


\begin{frame}{Overview}
	\setbeamertemplate{section in toc}[sections numbered]
	\tableofcontents[hideallsubsections]
\end{frame}


% ################# ORGANISATION #################


\section{Organisation}

\subsection{Who are we?}
\begin{frame}{Who are we?}

	\begin{columns}
		\column{0.50\textwidth}
		Florian Kluge\\
		Florian.Kluge@mailbox.tu-dresden.de



		\column{0.50\textwidth}
		Moritz Schulz\\
		Moritz.Schulz2@mailbox.tu-dresden.de\\
		@schokotets auf Telegram

	\end{columns}

\end{frame}

\subsection{What are we doing here?}
\begin{frame}{What are we doing here?}
	What are we doing here?
	\begin{itemize}
		\item Introduction to programming
		\item Getting to know the basics of Java
		\item Preparation for upcoming courses (e.g 'Softwaretechnologie', 2nd Semester)
	\end{itemize}
\end{frame}

\subsection{Structure}
\begin{frame}{Structure}
	\begin{itemize}[<+->]
		\item ~14 lessons
		\item Thursday, 13:00 - 14:30
		\item APB/006 (right here)
		\item 3G-Rule
		\item Attendance list
	\end{itemize}
\end{frame}

\begin{frame}{Structure}
	\begin{itemize}[<+->]
		\item This course is *voluntarily*
		\item Bored and want to leave? No problem!
		\item[]	.. but please contact us so we can invite students from the waiting list
		\item If you don't attend the course for two weeks in a row without notice we will give your slot to other students
	\end{itemize}
\end{frame}

% ################# Why Java? #################

\section{Why Java?}
\begin{frame}{Why Java?}
	\begin{itemize}
		\item Widely used programming language
		\item Introduction to object oriented programming (OOP)
		\item Platform-independent
		\item ... and much more
	\end{itemize}
\end{frame}

\subsection{Use cases}
\begin{frame}{Use cases}
	\begin{itemize}
		\item Android development
		\item Web applications
		\item Desktop GUI applications
		\item ... and much more
	\end{itemize}
\end{frame}


\subsection{Experience}
\begin{frame}{Do you have any programming experience already?}
	\begin{center}
		Do you have any programming experience already?\\
		\url{https://trivo25.github.io/tud-java-course/poll.html}\\
		or\\
		\url{https://strawpoll.com/6uh45fcvx}
	\end{center}
\end{frame}


% ################# Time to get started! ..almost #################

\section{Time to get started! ..almost}

\begin{frame}[fragile]{Time to get started! ..almost}
	\begin{center}
		Java OpenJDK 11 \url{https://adoptium.net/}
		\\
		Did you install it correctly? Time to find out!
	\end{center}

	\begin{lstlisting}[language=bash]
		$ javac --version
		> javac 11.0.12
	\end{lstlisting}
\end{frame}



\begin{frame}[fragile]{Time to get started! ..almost}
	\begin{center}
		Doesn't work? :(
		Use an online compiler!

		\url{https://www.jdoodle.com/online-java-compiler/}
	\end{center}
\end{frame}



% ################# lets go #################

\section{Let's go!}

\subsection{Your first task}
\begin{frame}[fragile]{Your first task}
	\begin{itemize}[<+->]
		\item Create a new folder
		\item Open the terminal and navigate into that folder using \lstinline[columns=fixed]{$ cd /to/my/folder}
		\item Create a new file by either typing\\ \lstinline[columns=fixed]{$ touch helloWorld.java}\\
		      or right-clicking in your folder \lstinline[columns=fixed]{Right click -> new -> text document}\\
		      and save it as a \lstinline[columns=fixed]{.java} file
	\end{itemize}
\end{frame}

\subsection{Your first task}
\begin{frame}[fragile]{Your first task}
	\begin{itemize}[<+->]
		\item now its time to write your first piece of code!
		\item[] \lstinputlisting[language=java]{../code_samples/HelloWorld.java}
	\end{itemize}
\end{frame}

\subsection{How to execute a java program}
\begin{frame}[fragile]{How to execute a java program}
	what we have to do now..
	\begin{itemize}[<+->]
		\item telling \lstinline[columns=fixed]{javac} to compile our \lstinline[columns=fixed]{helloWorld.java} file into a \lstinline[columns=fixed]{helloWorld.class}
		\item \lstinline[columns=fixed]{.class} files are 'bytecode' for the Java Virtual Machine (JVM)
		\item with \lstinline[columns=fixed]{$ java helloWorld} we can finally execute our first program!
	\end{itemize}
\end{frame}

\subsection{How to execute a java program}
\begin{frame}[fragile]{How to execute a java program}
	\begin{lstlisting}[language=bash]
		$ java helloWorld
		> Hello World!
	\end{lstlisting}
\end{frame}


\subsection{Time to play around}
\begin{frame}[fragile]{Time to play around}
	your next task
	\begin{itemize}[<+->]
		\item change the text you want to print in the \lstinline[columns=fixed]{helloWorld.java} file
		\item re-compile it into a \lstinline[columns=fixed]{.class} file and execute it again!
	\end{itemize}
\end{frame}


\subsection{What are we actually doing?}
\begin{frame}[fragile]{What are we actually doing?}
	\begin{itemize}[<+->]
		\item we are telling the computer what do to
		\item we list instructions for the computer
	\end{itemize}
\end{frame}


\subsection{Task numero 2!}
\begin{frame}[fragile]{Task numero 2!}
	\begin{center}
		Let's add a variable of type \lstinline[language=Java]{String}!
	\end{center}
\end{frame}

\subsection{Task numero 2!}
\begin{frame}[fragile]{Task numero 2!}
	\begin{center}
		\lstinputlisting[language=java]{../code_samples/VariableString.java}
	\end{center}
\end{frame}


\subsection{Task numero 2!}
\begin{frame}[fragile]{Task numero 2!}
	\begin{itemize}[<+->]
		\item We can re-use variables
		\item We can store data in them
	\end{itemize}
\end{frame}


\subsection{Task numero 3!}
\begin{frame}[fragile]{Task numero 3!}
	\begin{center}
		\lstinputlisting[language=java]{../code_samples/VariableStringName.java}
	\end{center}
\end{frame}


\subsection{Task numero 4!}
\begin{frame}[fragile]{Task numero 3!}
	\begin{itemize}[<+->]
		\item Besides \lstinline[language=Java]{String}s we also have variables of type \lstinline[language=Java]{int}\\
		\item \lstinline[language=Java]{int} represent whole numbers, like \lstinline[language=Java]{1, 18, 1337 or 420360}\\
		\item We can calculate \lstinline[language=Java]{int} with operators like \lstinline[language=Java]{+}, \lstinline[language=Java]{-}, \lstinline[language=Java]{*} and many more
	\end{itemize}
\end{frame}


\subsection{Task numero 4!}
\begin{frame}[fragile]{Task numero 3!}
	We now can..
	\begin{itemize}[<+->]
		\item Display text in the console \lstinline[language=Java]{System.out.println("Hello word!");}
		\item Declare variables like \lstinline[language=Java]{int} or \lstinline[language=Java]{String}
		\item Read input from the conole
		\item .. and know operators like \lstinline[language=Java]{+, - or *}
		\item[] Okay, what now?
		\item[] Let's build a calculator!
	\end{itemize}
\end{frame}
% ################# ending first lesson #################

\section{That's it (at least for today)}

\subsection{}
\begin{frame}{What will we do next lesson?}
	\begin{itemize}
		\item Deep dive into (more) variables and their operators
		\item Introducing functions and control flow
		\item and build more cool things!
	\end{itemize}
\end{frame}

\subsection{}
\begin{frame}{Links and resources}
	\begin{center}
		\url{https://trivo25.github.io/tud-java-course/}
	\end{center}

	\begin{figure}[htbp]
		\centerline{\includegraphics{assets/qr.png}}
		\label{fig}
	\end{figure}
\end{frame}

% nothing to do from here on
\end{document}
